\documentclass[aspectratio=169,11pt]{beamer}
\usepackage{../../../shared/theme/esmad_beamer_theme}
\usepackage{amsmath,amssymb}
\usepackage{booktabs}

\title[Policy Learning]{Policy Learning and Treatment Targeting}
\subtitle{From Causal Effects to Decisions}
\date{\today}

\newcommand{\E}{\mathbb{E}}

\begin{document}

\begin{frame}
\titlepage
\end{frame}

\begin{frame}{Learning Goals}
\begin{itemize}
\item distinguish treatment-effect estimation from policy learning;
\item define treatment rules, policy value, and regret;
\item translate CATE estimates into decisions under costs and thresholds;
\item evaluate policies using experimental or observational data;
\item understand targeting curves and budget-constrained assignment;
\item recognize uncertainty, fairness, and deployment risks.
\end{itemize}
\end{frame}

\begin{frame}{Outline}
\tableofcontents
\end{frame}

\section{From Effects to Decisions}

\begin{frame}{Treatment Rules}
After estimating
\[
\tau(x)=\E[Y(1)-Y(0)\mid X=x],
\]
we may choose an action
\[
\pi(x)\in\{0,1\}.
\]

\begin{block}{Decision question}
Which units should receive treatment under the objective and constraints we actually care about?
\end{block}
\end{frame}

\begin{frame}{Effect Estimation Is Not Policy Learning}
A CATE model estimates heterogeneous benefit.

A policy must additionally account for:
\begin{itemize}
\item treatment cost;
\item capacity;
\item harm;
\item uncertainty;
\item fairness;
\item operational feasibility.
\end{itemize}

\begin{alertblock}{Key distinction}
High predicted treatment effect is not automatically a treatment recommendation.
\end{alertblock}
\end{frame}

\section{Policy Value and Regret}

\begin{frame}{Policy Value}
For a rule \(\pi\),
\[
V(\pi)=\E[Y(\pi(X))].
\]

Policy learning seeks a rule with high expected value, not merely a model with low prediction error.
\end{frame}

\begin{frame}{Regret}
Relative to an optimal feasible rule \(\pi^\star\),
\[
R(\pi)=V(\pi^\star)-V(\pi).
\]

\begin{block}{Interpretation}
Regret is the opportunity cost of using a suboptimal assignment policy.
\end{block}
\end{frame}

\section{Costs and Constraints}

\begin{frame}{Net Benefit}
If treatment has cost \(c(X)\), the relevant quantity can be
\[
\tau(X)-c(X).
\]

A simple rule is
\[
\pi(x)=\mathbb{1}\{\widehat\tau(x)>c(x)\}.
\]
\end{frame}

\begin{frame}{Budget-Constrained Targeting}
Suppose only a fraction \(q\) of units can be treated.

A common strategy:
\begin{enumerate}
\item estimate treatment benefit;
\item rank units by expected net benefit;
\item treat the top \(q\) fraction;
\item validate policy value out of sample.
\end{enumerate}
\end{frame}

\begin{frame}{Capacity Changes the Target}
With scarce treatment capacity, the estimand is no longer simply an ATE.

The target becomes the expected value of a feasible assignment rule.

\begin{alertblock}{Implication}
A modest average effect can still support a valuable policy if treatment can be targeted well.
\end{alertblock}
\end{frame}

\section{Off-Policy Evaluation}

\begin{frame}{The Evaluation Problem}
We often want to evaluate a new rule using data generated by a different treatment policy.

Observed outcomes correspond to the action actually taken, not necessarily the action recommended by the new rule.
\end{frame}

\begin{frame}{Inverse-Propensity Policy Evaluation}
A simple estimator is
\[
\widehat V_{IPW}(\pi)
=
\frac{1}{n}\sum_i
\frac{\mathbb{1}\{D_i=\pi(X_i)\}Y_i}
{P(D_i=\pi(X_i)\mid X_i)}.
\]

\begin{alertblock}{Support requirement}
If the proposed policy recommends actions rarely observed in the data, evaluation becomes unstable.
\end{alertblock}
\end{frame}

\begin{frame}{Doubly Robust Policy Evaluation}
Doubly robust evaluation combines:
\begin{itemize}
\item an outcome model;
\item a propensity model;
\item a correction for observed actions matching the policy.
\end{itemize}

It protects against one nuisance-model failure under standard conditions, but not against invalid causal identification.
\end{frame}

\section{Targeting Curves}

\begin{frame}{Uplift vs Response Prediction}
Response prediction estimates
\[
\E[Y\mid X].
\]

Uplift or causal targeting estimates
\[
\E[Y(1)-Y(0)\mid X].
\]

\begin{alertblock}{Different objective}
High response does not imply high incremental treatment benefit.
\end{alertblock}
\end{frame}

\begin{frame}{Targeting Curves}
Rank units by predicted treatment effect and inspect cumulative policy value as the treated fraction grows.

\begin{block}{Use}
Targeting curves show whether the model concentrates causal benefit among highly ranked units.
\end{block}
\end{frame}

\begin{frame}{Qini-Style Comparison}
A Qini-style curve compares targeted treatment with random assignment under the same treatment budget.

The focus is decision value, not pointwise CATE accuracy.
\end{frame}

\section{Uncertainty and Validation}

\begin{frame}{Noisy CATEs Create Risky Policies}
Policies can over-target statistical noise when:
\begin{itemize}
\item samples are small;
\item overlap is weak;
\item effects are modest;
\item the same data are used for learning and evaluation.
\end{itemize}
\end{frame}

\begin{frame}{Honest Evaluation}
Use held-out data, cross-fitting, or a randomized validation experiment when possible.

\begin{itemize}
\item learn the rule on training data;
\item estimate value on separate data;
\item compare with simple baseline policies;
\item report uncertainty in policy value.
\end{itemize}
\end{frame}

\section{Fairness and Harm}

\begin{frame}{Targeting Can Create Harm}
A learned rule may:
\begin{itemize}
\item systematically exclude high-uncertainty groups;
\item reproduce historical allocation bias;
\item optimize short-term value at long-term social cost;
\item concentrate treatment benefits unevenly.
\end{itemize}
\end{frame}

\begin{frame}{Constrained Policy Learning}
Constraints may include:
\begin{itemize}
\item budget limits;
\item minimum coverage;
\item maximum harm probability;
\item group-level treatment-rate constraints;
\item interpretability requirements.
\end{itemize}

These belong inside the optimization problem, not as an afterthought.
\end{frame}

\section{Deployment}

\begin{frame}{Exploration}
Historical data may contain little information about actions rarely taken.

Deployment may therefore require:
\begin{itemize}
\item randomized exploration;
\item phased rollout;
\item guardrails;
\item stopping rules;
\item continuous monitoring.
\end{itemize}
\end{frame}

\begin{frame}{Policy Drift}
Policy value can change when populations, competitors, treatment delivery, or user behavior change.

A learned policy is therefore a monitored decision system, not a one-time model.
\end{frame}

\section{Applied Workflow}

\begin{frame}{Example: Regional Pricing Policy}
Suppose a pricing intervention has heterogeneous effects.

A policy-learning analysis asks:
\begin{itemize}
\item which regions have positive expected net benefit?
\item how many regions can be treated?
\item does targeting outperform random assignment?
\item how uncertain are the marginal decisions?
\item are some groups systematically excluded?
\end{itemize}
\end{frame}

\begin{frame}{A Credible Policy-Learning Workflow}
\begin{enumerate}
\item Define the action space and target population.
\item State the causal estimand and decision objective.
\item Estimate treatment heterogeneity under valid identification.
\item Define costs, harms, and constraints.
\item Learn the policy on training data.
\item Evaluate policy value on held-out or experimental data.
\item Compare against simple baselines.
\item Report uncertainty and regret.
\item Monitor drift after deployment.
\end{enumerate}
\end{frame}

\begin{frame}{Common Failure Modes}
\begin{itemize}
\item treating high predicted outcome as high treatment benefit;
\item learning and evaluating on the same data;
\item ignoring treatment cost;
\item ignoring weak overlap;
\item optimizing CATE accuracy instead of policy value;
\item deploying without monitoring;
\item omitting fairness and harm constraints.
\end{itemize}
\end{frame}

\begin{frame}{Synthesis: Decisions Need Their Own Target}
\begin{itemize}
\item Policy learning turns causal estimates into treatment rules.
\item Policy value, not effect-model accuracy, is the decision target.
\item Costs and constraints determine who should actually be treated.
\item Off-policy evaluation requires support for the actions recommended.
\item Targeting curves summarize incremental value under limited capacity.
\item Honest validation is essential because CATE estimates are noisy.
\item Fairness, harm, and monitoring belong inside the policy problem.
\end{itemize}
\end{frame}

\begin{frame}{Selected References}
\small
\begin{itemize}
\item Manski, C. F. (2004). Statistical Treatment Rules for Heterogeneous Populations. \textit{Econometrica}, 72(4), 1221--1246.
\item Qian, M., and Murphy, S. A. (2011). Performance Guarantees for Individualized Treatment Rules. \textit{Annals of Statistics}, 39(2), 1180--1210.
\item Kitagawa, T., and Tetenov, A. (2018). Who Should Be Treated? Empirical Welfare Maximization Methods. \textit{Econometrica}, 86(2), 591--616.
\item Athey, S., and Wager, S. (2021). Policy Learning with Observational Data. \textit{Econometrica}, 89(1), 133--161.
\item Zhao, Y.-Q., Zeng, D., Rush, A. J., and Kosorok, M. R. (2012). Estimating Individualized Treatment Rules Using Outcome Weighted Learning. \textit{JASA}, 107(499), 1106--1118.
\end{itemize}
\end{frame}

\contactslide

\end{document}
